\section{Odd-prime first jets at all higher levels}
\label{odd:section}

The residue formula loses both state directions on a supersingular fiber.
To retain their first vertical terms, we compare the original divided
forms on a common square-zero thickening. The calculation keeps the
deformation inside each tame block as well as the rescaling between
blocks. Neither a Hasse denominator nor a simple-root assumption occurs.
This differs from working in a unit-root quotient after making the
Hasse--Witt matrix invertible, as in \cite{BeukersVlasenko2021}; the
higher-Hasse framework discussed in \cite{Vlasenko2024} also has its own
lower-level invertibility hypotheses. The first-jet identity needed here
will follow directly from the original matrix.

\subsection{The common thickening and the statement}

In this section $p$ is odd, $p\nmid m$, and the residue field $k$ is
finite. Use the coefficient rings and parameters of the main statements,
and put
\[
 N=p^a,\qquad \sigma=1+p+\cdots+p^{a-1},\qquad
 \ell=\sigma-1,\qquad T=\bar t_a^m,
 \qquad J=I_{m,\eta}(x,y;\bar t_a).
\]
Write $H=H_p(T,J;\varepsilon_m)$, where
$\varepsilon_m=(-1)^{m+1}$. The cyclotomic valuation is
$v_{\pi_i}(p)=p^{i-1}(p-1)\ge2$. Thus
$\mathcal O_i/(\pi_i^2)$ has characteristic $p$ and a square-zero
maximal ideal of length one. For $k=\mathbb F_{p^f}$ the roots of
$X^{p^f}-X$ lift uniquely to this quotient, since its derivative is
$-1$. They form a coefficient field: sums and products remain roots,
and reduction identifies their field with $k$. Every element therefore
has a unique expression $c_0+c_1\pi_i$, with $c_0,c_1\in k$.

Identify the two quotients for $i=1,a$ by
\begin{equation}
 B=k[\epsilon]/(\epsilon^2),\qquad
 \pi_1\longmapsto\epsilon,\quad \pi_a\longmapsto\epsilon,
 \qquad s_1,s_a\longmapsto\eta(1+\epsilon),
 \label{odd:common-ring}
\end{equation}
where the displayed arrows mean the two specified maps to $B$.
Choose $t_1\in\mathcal O_1^\times$ whose image equals the entire image
$t$ of $t_a$ in $B$, not just its residue. The quotient map is
surjective, so such a choice exists. The four last center coordinates
$1,t,t,s$ and the eight successive coordinate substitutions in
\eqref{geom:charts} are then identical over $B$. They identify the two
original open models with one smooth surface $\mathcal U_B$.
This uses the actual blowup charts, not a general assertion that
blowup commutes with every nonflat base change.

The comparison in \eqref{odd:common-ring} is not the natural cyclotomic
embedding. For $a\ge2$, that embedding sends $\pi_1$ to
$\zeta_{p^a}^{p^{a-1}}-1$, whose image in $B$ is zero, not $\epsilon$.
The superscript $[2]$ below always uses the specified common quotient
and the matched time. Formula~\eqref{trace:integral-insertion} shows on
the torus that higher-order choices of the time lift do not affect
the resulting form over $B$; uniqueness on the terminal charts follows
from Lemma~\ref{trace:chart-injection}.

\begin{theorem}[Complete odd-prime first-jet factorization]
\label{odd:factorisation}
For every odd $p$, every $m$ prime to $p$, and every $a\ge2$ in the
finite-residue-field setting above, the entire original open surface
satisfies
\begin{equation}
 \alpha_{mp^a}^{[2]}
 =H^{\langle\sigma-1\rangle}\alpha_{mp}^{[2]}.
 \label{odd:factorisation-equation}
\end{equation}
Here $H^{\langle\sigma-1\rangle}$ is the common power of any local
regular lift of $H$. These powers glue even when $J$ has no global
lift. The equality includes all four terminal affine lines and does
not require a smooth energy fiber, ordinarity, or a simple Hasse root.
\end{theorem}

We prove the theorem through block identities that also supply the
characteristic-independent inputs to Section~\ref{two:section}.
In those algebraic identities the field containing $\eta$ need not be
finite, and a prime may be two whenever its square-zero quotient has
that characteristic. The full model comparisons and geometric ideal
statements retain their stated residue-field hypotheses.

\subsection{Integral blocks and the full internal deformation}

For the original matrix \eqref{main:matrix}, set
\[
 \mathsf B_s(z)=A(s^{m-1}z;t)\cdots A(z;t),\qquad r=s^m.
\]
If $s$ has exact order $mN$, then $r$ has exact order $N$, and the
original descending product is
$\mathsf B_s(r^{N-1}z)\cdots\mathsf B_s(z)$.

\begin{lemma}[Block insertion]
\label{odd:block-insertion}
Before reduction, for every prime $p$ and $N=p^a$,
\begin{equation}
 \alpha_{mN}=[z^{mN}]\operatorname{tr}\bigl(
 d\mathsf B_s(z)\mathsf B_s(r^{N-1}z)\cdots\mathsf B_s(rz)\bigr).
 \label{odd:insertion-equation}
\end{equation}
The right side is integral and includes all $m$ insertions inside
$d\mathsf B_s$; no additional factor $m$ is present.
\end{lemma}

\begin{proof}
Apply the cyclic insertion argument of Lemma~\ref{trace:insertion}
to the $N$ blocks. After rotating the differentiated block to the
first position, substitute $w=r^jz$. The remaining block indices
are $N-1,\ldots,1$, using $r^N=1$, and the extracted coefficient
is unchanged because $r^{jmN}=1$. Thus the $N$ block insertions
give $dI_{mN}=N$ times the right side of
\eqref{odd:insertion-equation}. Division occurs in characteristic
zero and leaves the displayed integral expression.
\end{proof}

Now work on $R=B[x^{\pm1},y^{\pm1}]$, keeping the actual
$t\in B^\times$. For any polynomial with scalar or one-form
coefficients define
\begin{equation}
 \mathcal P_m\left(\sum_i b_i z^i\right)=\sum_{j\ge0}b_{mj}Z^j.
 \label{odd:projection-definition}
\end{equation}
This is a coefficient projection, not a ring homomorphism. Its
valid multiplicative rule is
\begin{equation}
 \mathcal P_m\bigl(Q(z^m)V(z)\bigr)=Q(Z)\mathcal P_m V(z).
 \label{odd:projection-rule}
\end{equation}

Put $\mathsf A_j(z)=A(\eta^jz;t)$. Expanding every factor, including
those inside the tame block, gives
\begin{equation}
 \begin{split}
 \mathsf B_s&=\mathsf B_\eta(z;t)+\epsilon C(z),\\
 C(z)&=\sum_{j=0}^{m-1}j\,
 \mathsf A_{m-1}\cdots\mathsf A_{j+1}
 (z\partial_z\mathsf A_j)
 \mathsf A_{j-1}\cdots\mathsf A_0.
 \end{split}
 \label{odd:block-variation}
\end{equation}
Empty products are identity matrices. Although $C$ need only be
specified modulo $\epsilon$, the time in $\mathsf B_\eta(z;t)$
has not been replaced by $\bar t$. In particular, time derivatives
of state differentials remain in this formula.

\begin{lemma}[The deformed block trace]
\label{odd:block-trace}
Write $S_s=\operatorname{tr}\mathsf B_s$ and
$D_s=\det\mathsf B_s$. Then
\begin{equation}
 \begin{gathered}
 S_s(z)=S_*(z^m)+\epsilon V(z),\qquad \mathcal P_m V=0,\\
 S_*(Z)=T_*+j_*Z+c_sZ^2,\qquad
 T_*=t^m,\quad j_*=[z^m]S_s,\quad c_s=s^{m(m-1)},\\
 D_s(z)=d_s z^{3m},\qquad
 d_s=s^{3m(m-1)/2},\qquad
 dS_*=Z\,dj_*,\quad dD_s=0.
 \end{gathered}
 \label{odd:deformed-trace}
\end{equation}
The residue of $j_*$ is $J$. For odd $p$ the two state-constant
leading coefficients are
\begin{equation}
 c_s=1+\epsilon m(m-1),\qquad
 d_s=\varepsilon_m\left(1+\epsilon\frac{3m(m-1)}2\right).
 \label{odd:leading-jets}
\end{equation}
The integer $3m(m-1)/2$ in these formulas is formed before reduction.
\end{lemma}

\begin{proof}
At $s=\eta$ the cyclic trace is invariant under $z\mapsto\eta z$.
All nonmultiples of $m$ vanish, since $\eta^i-1$ is a unit for
$m\nmid i$. At $z=0$, the identities
$\operatorname{tr}A(0;t)=t$ and $\det A(0;t)=0$ give
$A(0;t)^m=t^{m-1}A(0;t)$. The constant trace coefficient is $t^m$.
The highest block coefficient is $s^{m(m-1)}E$, because the
highest coefficient of every original factor is a scalar multiple
of $E=\operatorname{diag}(1,0)$ and $E^m=E$. Multiplying the original
determinants gives $D_s=d_s z^{3m}$. These arguments remain valid
over $B$ with its actual unit time. Combining them with
\eqref{odd:block-variation} proves \eqref{odd:deformed-trace}.
Finally, $\eta^{m(m-1)}=1$ and
$\eta^{3m(m-1)/2}=\varepsilon_m$: for odd $m$ the exponent is a
multiple of $m$, and for even $m$ one uses $\eta^{m/2}=-1$.
Expanding $s=\eta(1+\epsilon)$ proves
\eqref{odd:leading-jets}.
\end{proof}

\begin{lemma}[Square-zero nonresonant terms]
\label{odd:projection}
Under \eqref{odd:deformed-trace}, put $D_*(Z)=d_sZ^3$.
For any polynomial $\Phi(S,D)$ over the coefficient ring,
\begin{equation}
 \mathcal P_m\bigl(\Phi(S_s,D_s)dS_s\bigr)
 =\Phi(S_*,D_*)dS_*.
 \label{odd:projection-equation}
\end{equation}
\end{lemma}

\begin{proof}
The square-zero expansion differs from the proposed right side
only by the projections of
\[
 \epsilon\Phi_S(S_*(z^m),D_*(z^m))V\,dS_*(z^m),
 \qquad
 \epsilon\Phi(S_*(z^m),D_*(z^m))dV.
\]
Both are zero by \eqref{odd:projection-rule}, since
$\mathcal P_mV=\mathcal P_m(dV)=0$. Thus no nonresonant term
contributes at this order. The resonant changes in $j_*,c_s,d_s$
have not been discarded.
\end{proof}

Return to odd $p$ and define
\[
 F_N=[z^{mN}]\operatorname{tr}
       (d\mathsf B_s\mathsf B_s^{N-1}),\qquad
 f_*=S_*^2-4D_*,\qquad
 H_*=[Z^{p-1}]f_*^{(p-1)/2}.
\]
The universal polynomial identity for $U_{N-1}$ proved in
Lemma~\ref{trace:residue} holds over nonreduced rings as well.
Equation~\eqref{trace:rank-two}, $dD_s=0$, and
Lemma~\ref{odd:projection} therefore give
\begin{equation}
 F_N=[Z^{N-1}]f_*^{(N-1)/2}\,dj_*.
 \label{odd:internal-coefficient}
\end{equation}
For completeness, the coefficient selection here includes all the
deformed scalar coefficients. Put $Q_*=f_*^{(p-1)/2}$; its support
lies in $[0,2p-2]$, because $\deg f_*\le4$, and
\[
 f_*^{(N-1)/2}=\prod_{i=0}^{a-1}Q_*^{p^i}.
\]
A contribution to $Z^{p^a-1}$ has lowest index congruent to $p-1$
modulo $p$. In $[0,2p-2]$ only $p-1$ has that congruence.
Subtracting it and dividing the remaining degree by $p$ repeats
the argument at every index. Consequently
\begin{equation}
 F_N=H_*^\sigma dj_*,\qquad F_p=H_*dj_*.
 \label{odd:internal-factorisation}
\end{equation}
Frobenius is applied only to scalar coefficients, never to a one-form.
The functions $T_*,j_*,H_*$ and the form $dj_*$ still contain the
full time jet. Since $\bar H_*=H$ and $p\mid\ell$ when $a\ge2$,
the power $H_*^\ell$ equals the power of any other lift of $H$.

\subsection{The weighted commutator and the external terms}

\begin{lemma}[Weighted rank-two identity]
\label{odd:commutator}
Let $M,X$ be two-by-two matrices over a commutative ring of
characteristic $p$, and let $N=p^a$. Then
\begin{equation}
 \sum_{j=1}^{N-1}jM^{N-1-j}XM^{j-1}
 =\operatorname{ad}_M^{N-2}(X),
 \qquad \operatorname{ad}_M(X)=[M,X].
 \label{odd:weighted-operator}
\end{equation}
For odd $p$, writing $S=\operatorname{tr}M$ and $D=\det M$ gives
\begin{equation}
 \sum_{j=1}^{N-1}jM^{N-1-j}XM^{j-1}
 =(S^2-4D)^{(N-3)/2}[M,X].
 \label{odd:weighted-commutator}
\end{equation}
No inverse of a determinant or discriminant is used.
\end{lemma}

\begin{proof}
In $\mathbb F_p[U,V]$ the geometric sum is
\[
 \sum_{j=0}^{N-1}U^{N-1-j}V^j=(U-V)^{N-1},
\]
by $U^N-V^N=(U-V)^N$. This equality is first proved in the
polynomial domain and then specialized, so cancellation is not
being performed in the matrix coefficient ring. Differentiate
with respect to $V$; the right side has coefficient $-(N-1)=1$.
Substituting the commuting operators of left and right
multiplication by $M$ proves \eqref{odd:weighted-operator}.

Over any commutative ring, Cayley--Hamilton gives
$M^2=SM-D\,\mathrm{id}$ and
$M^3=(S^2-D)M-SD\,\mathrm{id}$. Substitution in the cubic
commutator yields
\[
 \begin{split}
 \operatorname{ad}_M^3(X)
 &=M^3X-3M^2XM+3MXM^2-XM^3\\
 &=(S^2-D)[M,X]-3D[M,X]
 =(S^2-4D)[M,X].
 \end{split}
\]
For odd $N$, iterate this identity $ (N-3)/2 $ times to obtain
\eqref{odd:weighted-commutator}. It remains valid for $N=3$, with
exponent zero, and for scalar matrices or vanishing discriminant.
\end{proof}

Let $\mathsf B_0=\mathsf B_\eta(z;\bar t)$ be the residue block and set
\begin{equation}
 \begin{split}
 \Gamma_m(z)&=\operatorname{tr}\bigl(
 d\mathsf B_0[\mathsf B_0,z\partial_z\mathsf B_0]\bigr),\\
 \gamma_m(Z)&=\mathcal P_m\Gamma_m(z),\qquad
 f(Z)=(T+JZ+Z^2)^2-4\varepsilon_m Z^3.
 \end{split}
 \label{odd:external-data}
\end{equation}
The external speed is $r=s^m=1+m\epsilon$, whence
\[
 \mathsf B_s(r^jz)=\mathsf B_s(z)
       +\epsilon mj\,z\partial_z\mathsf B_0(z).
\]
In the descending product of \eqref{odd:insertion-equation}, an
insertion at position $j$ has $N-1-j$ unchanged blocks on its
left and $j-1$ on its right. Thus
\eqref{odd:weighted-commutator} fixes the positive sign and gives
\begin{equation}
 \begin{split}
 \alpha_{mN}^{[2]}&=F_N+\epsilon m\beta_N,\\
 \beta_N&=[z^{mN}]\Gamma_m(z)f(z^m)^{(N-3)/2}
        =[Z^N]\gamma_m(Z)f(Z)^{(N-3)/2}.
 \end{split}
 \label{odd:external-coefficient}
\end{equation}
All other factors in an external correction are reduced because
that correction already contains $\epsilon$. This does not reduce
the internal term $F_N$ or its time jet.

\begin{lemma}[Actual block support]
\label{odd:support}
For the original residue block in any characteristic,
\[
 \operatorname{supp}_z\Gamma_m\subset[1,6m-2],\qquad
 \operatorname{supp}_Z\gamma_m\subset
 \begin{cases}[1,4],&m=1,\\{}[1,5],&m\ge2.
 \end{cases}
\]
\end{lemma}

\begin{proof}
The highest coefficient of $\mathsf B_0$ is the state-constant
matrix $E$, so $\deg d\mathsf B_0\le2m-1$. The degree $4m$
coefficient of $[\mathsf B_0,z\partial_z\mathsf B_0]$ is
$[E,2mE]=0$, while its constant term vanishes. Its degree is
therefore at most $4m-1$. Multiplication by $d\mathsf B_0$ and
projection to multiples of $m$ give the asserted bounds.
\end{proof}

The fifth possible resonant degree is essential for general $m$;
it cannot be removed by using the degree bound for a single original
matrix. Put
\[
 Q=f^{(p-1)/2},\qquad G=\gamma_m f^{(p-3)/2}.
\]
Lemma~\ref{odd:support} gives the actual bounds
\begin{equation}
 \operatorname{supp}Q\subset[0,2p-2],\qquad
 \operatorname{supp}G\subset[1,2p-1].
 \label{odd:digit-support}
\end{equation}
These also cover $p=3$, when $G=\gamma_m$ and its upper bound is
five. For $a\ge2$, split the exponent in
\eqref{odd:external-coefficient} to get
\[
 \beta_N=[Z^{p^a}]G(Z)\prod_{i=1}^{a-1}Q(Z)^{p^i}.
\]
Every potentially contributing degree satisfies
\[
 n_0+pn_1+\cdots+p^{a-1}n_{a-1}=p^a,
 \quad 1\le n_0\le2p-1,\quad
 0\le n_i\le2p-2\quad(i\ge1).
\]
Reduction modulo $p$ forces $n_0=p$, the only multiple of $p$
in its permitted interval. The remaining degree, after subtraction
and division by $p$, is $p^{a-1}-1$. At each remaining index
the only permissible value congruent to $p-1$ is $p-1$.
The coefficient extraction is consequently exact:
\begin{equation}
 \beta_N=H^{p+p^2+\cdots+p^{a-1}}\beta_p
        =H^\ell\beta_p.
 \label{odd:external-factorisation}
\end{equation}
Only $Q$, a scalar polynomial, is raised to Frobenius powers.

\subsection{Global identity and the two-coefficient ideal}

\begin{proof}[Completion of the proof of Theorem~\ref{odd:factorisation}]
The internal and external calculations give, on the common torus,
\[
 \alpha_{mN}^{[2]}=H_*^\sigma dj_*+\epsilon mH^\ell\beta_p,
 \qquad
 \alpha_{mp}^{[2]}=H_*dj_*+\epsilon m\beta_p.
\]
The two formulas use exactly the same $s,t,\mathsf B_s$ over $B$.
They therefore retain both the tame-block deformation and the
actual time term, including the change in $dj_*$.

If two local lifts of $H$ differ by $\epsilon b$, their $p$th
powers agree. Since $p\mid\ell$, their $\ell$th powers agree too.
These powers glue to a global regular $H^{\langle\ell\rangle}$,
and on the torus this function is $H_*^\ell$. Multiplying the
second displayed formula by it gives the first. The external
factor $m$ is identical at both levels and introduces no further
factor into the comparison.

Both divided forms are globally regular by
Proposition~\ref{trace:regular}. Their difference after multiplication
is therefore a regular one-form on $\mathcal U_B$, zero on the torus.
The four terminal chart rings are the appropriate localizations of
$B[u,v]$ by $1+uv,t+uv,t+uv,s+uv$, and their torus intersections
are obtained by inverting $u$. Lemma~\ref{trace:chart-injection}
makes restriction injective in the free cotangent module even
over this nonreduced base. The difference vanishes on each of
the four charts, proving the theorem on the whole open surface.
No regularity assertion for $\beta_p$ by itself is needed.
\end{proof}

\begin{corollary}[Exact odd-prime truncated ideal]
\label{odd:ideal}
At every point of an original complete smooth finite energy fiber,
for every $a\ge2$ and odd $p$, one has
\begin{equation}
 \mathfrak c(\alpha_{mp^a})+(\pi_a^2)
 =\bigl(\pi_a^2,\widetilde H^\sigma,
                 \pi_a\widetilde H^{\sigma-1}\bigr).
 \label{odd:ideal-equation}
\end{equation}
The lift $\widetilde H$ is arbitrary. If $h$ is a root of multiplicity
$e_h$ of $H_p(T,h;\varepsilon_m)$, then, after an allowed unramified
residue-field extension and completion, a coordinate $z$ with
$\bar z=J-h$ gives the ideal
\begin{equation}
 (\pi_a^2,z^{e_h\sigma},\pi_a z^{e_h(\sigma-1)}).
 \label{odd:multiple-root}
\end{equation}
\end{corollary}

\begin{proof}
For a scalar $g$ and a one-form $\omega$ in a free module,
$\mathfrak c(g\omega)=g\mathfrak c(\omega)$, by taking both
coefficients in a basis; $g$ need not be a unit. Reduce the complete
first-layer ideal of Theorem~\ref{first:ideal} to $B$ and multiply
its two generators by $H^{\langle\ell\rangle}$, using
Theorem~\ref{odd:factorisation}. Taking the inverse image under
reduction modulo $\pi_a^2$ gives \eqref{odd:ideal-equation}.
Thus the full-form comparison, rather than a tangential quotient
alone, supplies both generators.

At a root of multiplicity $e_h$, the residue Hasse function is
$(J-h)^{e_h}V(J)$ with $V(h)\ne0$. Smoothness permits a local
coordinate $z$ lifting $J-h$, and the lift of $V(J)$ is a unit.
Choose their product as the lift of $H$ in
\eqref{odd:ideal-equation}, and remove unit factors from the
generators. This proves \eqref{odd:multiple-root} without assuming
$e_h=1$. The permissible local completions or finite unramified
extensions do not alter the displayed original local identity.
\end{proof}

On an ordinary smooth fiber the ideals in
\eqref{odd:ideal-equation} are unit ideals. Along an allowed
unramified state lifting a supersingular point, $z$ has positive
$\pi_a$-order, so \eqref{odd:multiple-root} gives common order
at least two for the divided form. It gives neither exact order
two nor an untruncated higher-level critical ideal. The uniform
state-order summary is given at the end of the next section.
